\documentclass[a4paper, 12pt]{article}
\usepackage[left=3cm, top=3cm, right=2cm, bottom=2cm]{geometry}
%pacote que define as margens
\usepackage[utf8]{inputenc} %pacote que normaliza a acentuação
\usepackage[brazil]{babel} %pacote que define o idioma
\usepackage{graphicx} %pacote que permite a inserção de imagens
\pagestyle{myheadings} %pacote que define o cabeçalho com o número de página no canto superior direito
\usepackage{setspace} %pacote que permite especificar o espaçamento entre linhas no documento
\usepackage{indentfirst} %pacote que faz indentação na primeira linha do parágrafo
\setlength{\parindent}{1cm} %pacote que define o tamanho da indentação
\usepackage{float} %Pacote que permite tabelas flutuarem em qualquer posição
\usepackage[alf,abnt-repeated-title-omit=yes,abnt-emphasize=bf,abnt-etal-list=0]{abntex2cite} %define os parâmetros necessários para as referências e citações de acordo com a ABNT
\usepackage{amsthm}
\usepackage{amsmath}
\usepackage{amsfonts}
\usepackage{subcaption} % Pacote que permite que imagens possam ser inseridas em composições de imagens lado a lado com o \subfigure
\usepackage{amssymb} %4 Pacotes responsáveis pela formatação matemática
\usepackage{rotating} %permite que textos, figuras e tabelas possam ser rotacionados

\begin{document}
	
	\thispagestyle{empty}

\begin{center}
	\begin{figure}[h]
		\centering
		\includegraphics[width=0.20\linewidth]{pics/ufpa}
		\label{fig:ufpa}
	\end{figure}


	\vspace{1cm}
	\large \uppercase{UNIVERSIDADE FEDERAL DO PARÁ}\\
	\large \uppercase{INSTITUTO DE TECNOLOGIA}\\
	\large \uppercase{BACHARELADO EM ENGENHARIA MECÂNICA}\\
	\vspace{5cm}
	\large \uppercase{NÚCLEO AVANÇADO DE ANÁLISE TENSÕES}\\
	\vspace{1cm}
	\large \uppercase {SISTEMA DE AQUISIÇÃO REMOTA PARA EXTENSOMETRIA} \\
	\vspace{6cm}
	\large {BELÉM/PA \\ 2023}


	%FOLHA DE ROSTO
\end{center}
	\newpage
	\setcounter{page}{1} %comando que redefine a contagem de páginas
	\thispagestyle{empty}

\begin{center}
	LEONARDO DANTAS \\

	\vspace{5cm}

	NÚCLEO AVANÇADO DE ANÁLISE TENSÕES\\
	SISTEMA DE AQUISIÇÃO REMOTA PARA EXTENSOMETRIA\\
\end{center}
	\vspace{5cm}




\singlespacing
\hspace{8cm} % posicionando a caixa de texto
\begin{minipage}{7cm}
	\indent Artigo apresentado para o Núcleo Avançado de Análise de Tensões (NAAT), para o desenvolvimento de nova técnica de aquisição de dados.
\end{minipage}
\vspace{7cm}

\onehalfspacing
\begin{center}
	BELÉM/PA\\
	2023
\end{center}

\newpage
\thispagestyle{empty}
\begin{center}
	ALAN HENRIQUE PEREIRA MIRANDA - MATRICULA:202102140072 \\
	\vspace{5cm}

	Núcleo Avançado de Análise de Tensões\\
	SISTEMA DE AQUISIÇÃO REMOTA PARA EXTENSOMETRIA \\
	\vspace{3cm}

\end{center}
\singlespacing
\hspace{8cm} % posicionando a caixa de texto
\begin{minipage}{7cm}
	\indent Relatório, apresentado à Universidade da Amazônia, como parte das exigências para a obtenção de aprovação disciplinar.\\

	\vspace{2.5cm}

	Belém-PA 21 de maio de 2023 \\
	\vspace{3cm}
\end{minipage}

\onehalfspacing
\begin{center}

	EXAMINADOR\\
	\vspace{2cm}
	\rule{10cm}{0.15mm} \\
	Prof: Dr. Leonardo Dantas \\
	Universidade Federal do Pará - UFPA
\end{center}

	
	\section{Introdução}
\label{Introducao}

    \indent O circuito proposto é mostrado na Figura 1. Como parâmetros iniciais, determinou-se os valores abaixo.

\begin{figure}[h!]
    \centering
    \begin{circuitikz}
        \draw (0,0)
        (0,0) node[npn, tr circle](Q1){BC548}
        (Q1.base) node[anchor=east]{B}
        (Q1.collector) node[anchor=south]{C}
        (Q1.emitter) node[anchor=north]{E}
        (-2,0) -- (-1,0)
        (-2,0) to [R, l=10K{$\Omega$}] (-2,3.5)
        (0,0.9) to [R, l=3.9K{$\Omega$}] (0,3.5)
        (-2,3.5) -- (0,3.5)
        (-1, 3.5) -- (-1, 4) node[vcc]{$V_{cc}$}
        (-2,0) to [R, l=2.2K{$\Omega$}] (-2,-3.5)
        (-2,-3.5) -- (2, -3.5)
        (0, -0.9) to [eC, l=50<\mu\farad>] (0, -3.5)
        (0, -1) -- (2, -1)
        (2,-1) to[R, l=\SI{820}{$\Omega$}] (2,-3.5)
        (0,-3.5) -- (0,-4.5) node[tlground]{GND}
        (-2,0) to [eC, *-o, l=10<\mu\farad>] (-4,0) node[left]{IN}
        (0,0.9) to [eC, *-o, l=10<\mu\farad>] (2,0.9)  node[right]{OUT}
        ;
    \end{circuitikz}\label{fig:figure}
    \caption{Circuito proposto}

\end{figure}





	
	\section{Objetivos}
\label{sec:objetivos}

    \indent Este relatório pretende descrever experimentos realizados com um circuito amplificador de um único transistor. O circuito foi anteriormente projetado com auxílio do professor em sala, e os dados obtidos durante os testes serão usados para a montagem do circuito e validação.



	
	
	\addcontentsline{toc}{section}{Referências} %adiona a seção das referências no sumário
	\bibliographystyle{abntex2-alf}
	\bibliography{referencias}
	
\end{document}
